% Draft subsection -- "Grid Resolution Sensitivity and Its Impact on Vehicle
% Classification and Counting". All numbers: grid_resolution_analysis/
% (cross_validation_oof_summary.csv, final_numbers.json, timing_trials_summary.json).
\subsection{Grid Resolution Sensitivity and Its Impact on Vehicle Classification and Counting}
\label{sec:grid-sensitivity}

The lattice resolution fixes the cell size from which motion, object extent
and the span-based class label are derived. We therefore measured six
resolutions ($4{\times}6$, $6{\times}6$, $6{\times}8$, $8{\times}8$,
$8{\times}10$, $10{\times}12$, rows$\times$columns) on the two-lane serial
workload. The pipeline was otherwise unchanged.

\textbf{Protocol.} UA-DETRAC provides no camera-site labels. We grouped the 26
tuning videos into 10 sites by inspecting their static backgrounds. Grid choice
and the classification thresholds were then selected by leave-one-site-out
cross-validation on these videos only. In each fold, thresholds were tuned per
grid on the training sites, and a grid was chosen by a predefined objective. The
objective combines macro-F1, the weakest class F1, counting event-F1, detection
recall, fragmentation and unmatched detections, subject to a real-time
throughput guard. The resulting configuration was frozen before testing. Testing
used 44 further UA-DETRAC videos (67{,}206 frames). They had not been used
earlier in this work, and their scenes differ from all tuning and benchmark
sites. The 12 benchmark videos used elsewhere in this paper share sites with the
tuning data and had informed an earlier exploratory analysis, so their grid
results are retrospective only. Counting was evaluated against line-crossing
events extracted from the annotated trajectories. Events were matched one-to-one
within 1\,s and 80\,px. In the combined two-lane score, both lanes compete for
the same ground-truth events, so no vehicle can be credited twice.

\textbf{Classification.} With the original fixed thresholds, macro-F1 rose
monotonically with resolution (cross-validated 0.12 at $4{\times}6$ to 0.31 at
$10{\times}12$). Once thresholds were tuned per grid, the trend vanished
(cross-validated 0.30--0.35 for all grids). On the unseen-site test set, no
grid differed significantly from $6{\times}6$ (Table~\ref{tab:grid}; 95\%
session-level bootstrap). The apparent resolution effect was therefore mainly a
calibration effect of the fixed normalised-span thresholds. All tuned
classifiers exceeded the always-Car baseline on the test set (0.40--0.43 vs.\
0.31), but Van F1 remained at or below 0.16.

\textbf{Counting.} Finer grids detected slightly more vehicles (recall 0.883 at
$6{\times}6$ to 0.908 at $10{\times}12$) and produced more correct counts.
However, they also produced more false and duplicate counts. At $10{\times}12$
the system reported 1{,}728 crossings against 720 countable ground-truth
crossings, and the per-video absolute error rose from 14.9 ($6{\times}8$) to
23.1. Fragmentation grew from 1.5\% to 9.5\%. Diagnostic variants showed that
fragmentation arises mostly from single empty cell rows inside a vehicle, where
frame differencing leaves the interior unmarked; bridging such gaps removed
about 87--89\% of it. Diagonal-only adjacency accounts for a further 35--40\%.
Neither variant improved classification. Independently of resolution, the
counter registers only downward crossings, which excludes 44\% of the test
crossings.

\textbf{Selected configuration.} The predefined rule selected $6{\times}8$ with
tuned thresholds. It was the only grid whose objective improved significantly on
$6{\times}6$ on the test set, at 92\,FPS versus 97\,FPS on the same
two-core VM. The margin over $8{\times}8$ and $10{\times}12$ lies within
fold-to-fold variation. We therefore do not claim that any resolution is
universally best. We conclude that, beyond about $6{\times}8$, a finer
lattice trades counting reliability for marginal and statistically unsupported
classification gains. The span features and the single-direction counter limit
performance more than the grid does. All results come from one dataset with fixed
ROIs; external validation remains future work.

\begin{table}[t]
\centering
\caption{Grid resolution on 44 unseen-site UA-DETRAC test videos (frozen,
per-grid tuned thresholds; combined lanes). Count F1/MAE use downward
trajectory crossings. FPS: median of 3 serial trials on a 2-vCPU VM.}
\label{tab:grid}
\setlength{\tabcolsep}{3pt}
\begin{tabular}{lccccccc}
\hline
Grid & Macro-F1 & Recall & Frag. & Pred./GT & Count F1 & MAE & FPS\\
\hline
$6{\times}6$   & 0.405 & 0.883 & 1.5\% & 1377/720 & 0.314 & 16.0 & 97.2\\
$6{\times}8$*  & 0.416 & 0.892 & 2.3\% & 1318/720 & 0.336 & 14.9 & 91.9\\
$8{\times}8$   & 0.402 & 0.894 & 3.8\% & 1428/720 & 0.342 & 16.6 & 86.3\\
$10{\times}12$ & 0.430 & 0.908 & 9.5\% & 1728/720 & 0.337 & 23.1 & 81.4\\
\hline
\multicolumn{8}{l}{\footnotesize *selected by site-grouped cross-validation; always-Car macro-F1 = 0.31.}\\
\end{tabular}
\end{table}
